\documentclass[11pt]{article}
\usepackage[utf8]{inputenc}
\usepackage[T1]{fontenc}
\usepackage[a4paper,margin=1in]{geometry}
\usepackage{lmodern}
\usepackage{microtype}
\usepackage{amsmath,amssymb,mathtools}
\usepackage{mathrsfs}
\usepackage{tikz-cd}
\usepackage{enumitem}
\usepackage{longtable,booktabs,array}
\usepackage[hidelinks]{hyperref}
\setlength{\parindent}{0pt}
\setlength{\parskip}{0.55em}
\emergencystretch=3em
\hbadness=10000
\hfuzz=20pt
\newcommand{\K}{K}
\newcommand{\Gr}{\operatorname{Gr}}
\newcommand{\Fil}{\operatorname{Fil}}
\newcommand{\ch}{\operatorname{ch}}
\newcommand{\Todd}{\operatorname{Todd}}
\newcommand{\cl}{\operatorname{cl}}
\newcommand{\Parf}{\operatorname{Parf}}
\newcommand{\Tor}{\operatorname{Tor}}
\newcommand{\Spec}{\operatorname{Spec}}
\newcommand{\RRR}{\mathrm{RRR}}
\providecommand{\codim}{\operatorname{codim}}
\providecommand{\Inv}{\operatorname{Inv}}
\providecommand{\Pic}{\operatorname{Pic}}
\providecommand{\Z}{\mathbb Z}
\providecommand{\Q}{\mathbb Q}
\providecommand{\Ahat}{\widehat A}
\providecommand{\Ahp}{\widehat A^{+}}
\providecommand{\That}{\widehat T}
\providecommand{\rad}{\operatorname{rad}}
\providecommand{\Gl}{\operatorname{Gl}}
\providecommand{\GL}{\operatorname{GL}}
\providecommand{\ch}{\operatorname{ch}}
\providecommand{\Todd}{\operatorname{Todd}}
\providecommand{\Gr}{\operatorname{Gr}}

% Core macro additions
\providecommand{\Atil}{\widetilde A}
\providecommand{\Ctil}{\widetilde C}
\providecommand{\K}{\mathrm K}
\providecommand{\cC}{\mathcal C}
\providecommand{\cF}{\mathscr F}
\providecommand{\eps}{\varepsilon}
\providecommand{\ellc}{\ell}
\providecommand{\Tor}{\operatorname{Tor}}
\providecommand{\Coker}{\operatorname{Coker}}
\providecommand{\Ker}{\operatorname{Ker}}
\providecommand{\Supp}{\operatorname{Supp}}
\providecommand{\codim}{\operatorname{codim}}
\providecommand{\Gr}{\operatorname{G}}
\providecommand{\rank}{\operatorname{rank}}
\providecommand{\rang}{\operatorname{rang}}
\providecommand{\cO}{\mathcal O}

\hypersetup{hypertexnames=false}
\providecommand{\R}{\mathbb R}
\providecommand{\C}{\mathbb C}
\providecommand{\N}{\mathbb N}
\providecommand{\Zell}{\mathbb Z_\ell}
\providecommand{\Qell}{\mathbb Q_\ell}
\providecommand{\C}{\mathbb C}
\providecommand{\Ch}{\operatorname{Ch}}
\providecommand{\Sym}{\operatorname{Sym}}
\providecommand{\Td}{\operatorname{Td}}
\begin{document}

% Final-index helper macros
\providecommand{\Qcoh}{\operatorname{Qcoh}}
\providecommand{\Coh}{\operatorname{Coh}}
\providecommand{\Parf}{\operatorname{Parf}}
\providecommand{\Lib}{\operatorname{Lib}}
\providecommand{\Loclib}{\operatorname{Loclib}}
\providecommand{\Proj}{\operatorname{Proj}}
\providecommand{\Mod}{\operatorname{Mod}}
\providecommand{\Ob}{\operatorname{Ob}}
\providecommand{\Hom}{\operatorname{Hom}}
\providecommand{\toramp}{\operatorname{tor.amp}}
\providecommand{\parfamp}{\operatorname{par.amp}}
\providecommand{\coh}{\operatorname{coh}}
\providecommand{\qcoh}{\operatorname{qcoh}}
\providecommand{\parf}{\operatorname{parf}}
\providecommand{\torf}{\operatorname{torf}}
\providecommand{\amp}{\operatorname{amp}}
\providecommand{\num}{\operatorname{num}}
\providecommand{\naif}{\operatorname{naif}}
\providecommand{\alg}{\operatorname{alg}}
\providecommand{\Pic}{\operatorname{Pic}}
\providecommand{\Gr}{\operatorname{Gr}}
\providecommand{\Fil}{\operatorname{Fil}}
\providecommand{\Ch}{\operatorname{Ch}}
\providecommand{\Todd}{\operatorname{Todd}}
\providecommand{\Diff}{\operatorname{Diff}}
\providecommand{\Dcat}{\mathrm D}
\providecommand{\Tr}{\operatorname{Tr}}
\providecommand{\Z}{\mathbb Z}
\providecommand{\Q}{\mathbb Q}

\newpage
\section*{Index terminologique}
\begingroup\small
\begin{longtable}{@{}>{\raggedright\arraybackslash}p{0.64\linewidth}>{\raggedright\arraybackslash}p{0.14\linewidth}>{\raggedright\arraybackslash}p{0.12\linewidth}@{}}
\toprule
 & Exp. & §\\
\midrule
\endhead
additive (application) & I & 6.1\\
additive (application) & IV & 1.5\\
\(K-\lambda\)-algèbre augmentée & V & 3.9\\
algébriquement équivalent à zéro (dans \(K^\bullet\)) & X App. & 7.11\\
algébriquement équivalent à zéro (faisceau inversible) & XIII & 4.4\\
d'amplitude \(\mathcal C_0\)-parfaite contenue dans \([a,b]\) (complexe) & I & 4.7\\
d'amplitude \(\mathcal C_0\)-parfaite finie (complexe) & I & 4.7\\
d'amplitude parfaite contenue dans \([a,b]\) (complexe) & I & 4.8\\
d'amplitude parfaite finie (complexe) & I & 4.8\\
d'amplitude plate contenue dans \([a,b]\) (complexe) & I & 5.2\\
d'amplitude plate finie (complexe) & I & 5.2\\
d'amplitude parfaite relativement à \(f\) (ou \(Y\)) contenue dans \([a,b]\) (complexe) & III & 4.1\\
d'amplitude parfaite contenue dans \([-n,0]\) (morphisme) & III & 4.1\\
d'amplitude parfaite finie relativement à \(f\) (complexe) & III & 4.1\\
d'amplitude plate relativement à \(f\) (ou \(Y\)) contenue dans \([a,b]\) (complexe) & III & 3.1\\
d'amplitude plate contenue dans \([-n,0]\) (morphisme) & III & 3.1\\
d'amplitude plate finie relativement à \(f\) (complexe) & III & 3.1\\
amplitude projective & I & 4.12.1\\
\(\lambda\)-anneau & V & 2.4\\
\(\lambda\)-anneau engendré par générateurs et relations & V & 4.6\\
\(\lambda\)-anneau libre engendré par une famille de générateurs & V & 4.4\\
anneau de Chern & V & 6.2\\
application polynôme & X & 2.1.1\\
binomial (anneau) & V & 2.7\\
caractère de Chern & V & 6.4\\
caractère de Todd & V & 1.16\\
\(S\)-catégorie abélienne & I & 1.1.1\\
\(S\)-catégorie abélienne plate & I & 1.1.1\\
\(S\)-catégorie additive & I & 1.1.1\\
\(S\)-catégorie triangulée & I & 1.1.1\\
classe de Chern & V & 6.7\\
classe de Chern complétée & V & 6.7\\
classe de Chern totale & V & 6.7\\
codimension (d'une immersion régulière) & VII & 1.6\\
cohérent (complexe) & I & 3.1\\
cohéreur & II & 3.2\\
complexe cotangent relatif & VIII & 2.3\\
complexe elliptique & II App. II & 1.2.5\\
cup-produit & I & 8.2\\
\(k\)-cycle positif spécial & XIII & 6.9\\
degré (d'une application polynôme) & X & 2.1.1\\
degré (d'un faisceau inversible) & XIII & 2.2\\
degré (d'un morphisme de schémas) & X & 4.4\\
dense (sous-ensemble d'une catégorie abélienne ou triangulée) & IV & 1.7\\
dimension parfaite & I & 4.14.1\\
dimension relative virtuelle & VIII & 1.9\\
divisoriel & II & 2.2.5\\
dual & I & 7.3\\
dualité parfaite & I & 7.8\\
équivalence numérique (dans \(\operatorname{Gr}^\bullet(X)\) et \(\operatorname{Gr}_\bullet(X)\)) & X & 4.5\\
\(\tau\)-équivalence & X & 2.4.1\\
\(\tau\)-équivalence & XIII & 4.4\\
espace pseudo-analytique complexe & II & 2.4.7\\
espace pseudo-analytique complexe de Stein & II & 2.4.7\\
exact (sous-ensemble d'une catégorie abélienne ou triangulée) & IV & 1.7\\
\((b)\)-faisceau & XIII & 1.5\\
faisceau conormal & VII & 1.6\\
faisceau cotangent relatif & II App. II & 1.2.1\\
faisceau tangent relatif & II App. II & 1.2.1\\
\((b)\)-famille & XIII & 1.12\\
famille ample & II & 2.2.4\\
famille de classes de faisceaux cohérents & XIII & 1.12\\
famille de classes de faisceaux cohérents limitée & XIII & 1.12\\
famille de classes de faisceaux cohérents limitée par un faisceau cohérent & XIII & 1.12\\
fibré cotangent relatif & II App. II & 1.2.1\\
fibré tangent relatif & II App. II & 1.2.1\\
fidèlement plat (morphisme de topos annelés) & I & 2.16.2\\
filtration de \(\lambda\)-algèbre augmentée & V & 3.10\\
\(S\)-foncteur additif & I & 1.1.3\\
\(S\)-foncteur exact & I & 1.1.3\\
formule de projection & IV & 2.11\\
groupe de Grothendieck & IV & 2.2\\
\(\lambda\)-homomorphisme & V & 2.1\\
homomorphisme de spécialisation & X App. & 7.3\\
homomorphisme de spécialisation & X App. & 7.16\\
\(\lambda\)-idéal & V & 4.5\\
indice analytique & II App. II & 3.2.2\\
d'intersection complète (morphisme) & VIII & 1.1\\
\(n\)-isomorphisme & I & 1.4.3\\
libre homogène (module gradué) & VI & 1.2.2\\
localement stable par noyau d'épimorphisme & I & 1.2\\
localement quasi-relevable & I & 1.2\\
localement quasi-stable par noyau d'épimorphisme & I & 1.2\\
localement relevable & I & 1.2\\
nombre de Chern & X & 4.2\\
nombre d'intersection & X & 4.1\\
nombre d'intersection (pour une famille de faisceaux inversibles) & X & 4.4\\
numériquement équivalent à zéro (faisceau inversible) & XIII & 4.4\\
opérateur différentiel universel d'ordre \(n\) & II App. II & 1.2.1\\
opérations d'Adams & V & 7.1\\
\(\mathcal C_0\)-parfait (complexe) & I & 4.7\\
parfait (complexe) & I & 4.8\\
parfait (morphisme) & III & 4.1\\
parfait relativement à \(f\) (complexe) & III & 4.1\\
précohérent & I & 3.1\\
de \(n\)-\(\mathcal C_0\)-présentation finie & I & 2.8\\
de présentation finie & I & 2.8\\
d'\(\infty\)-\(\mathcal C_0\)-présentation finie & I & 2.8\\
\(n\)-\(\mathcal C_0\)-pseudo-cohérent (complexe) & I & 2.2\\
\(\mathcal C_0\)-pseudo-cohérent (complexe) & I & 2.2\\
\(n\)-pseudo-cohérent (complexe) & I & 2.15\\
\(n\)-pseudo-cohérent (morphisme) & III & 1.2\\
pseudo-cohérent (complexe) & I & 2.15\\
pseudo-cohérent (morphisme) & III & 1.2\\
\(n\)-pseudo-cohérent relativement à \(f\) (complexe) & III & 1.2\\
pseudo-cohérent relativement à \(f\) (complexe) & III & 1.2\\
\((b)\)-polynôme & XIII & 1.9\\
polynôme de Snapper & X & 2.3.1\\
polynôme universel & V & 1.15\\
pré-\(\lambda\)-anneau & V & 2.1\\
\(n\)-quasi-isomorphisme & I & 1.4.3\\
rang (d'un complexe parfait) & I & 6.11\\
\(m\)-régulier (faisceau) & XIII & 1.1\\
régulier (homomorphisme) & VII & 1.1\\
régulier (idéal) & VII & 1.4\\
régulière (immersion) & VII & 1.4\\
sous-\(S\)-catégorie triangulée & I & 1.1.4\\
spécial (\(\lambda\)-anneau) & V & 2.4\\
spécial (morphisme) & III & 1.1.4\\
de Stein (topos localement annelé) & II & 2.4.1\\
strictement \(\mathcal C_0\)-parfait & I & 2.1\\
strictement \(\mathcal C_0\)-\(n\)-pseudo-cohérent & I & 2.1\\
strictement \(\mathcal C_0\)-pseudo-cohérent & I & 2.1\\
système minimum de générateurs & VII & 1.4.1\\
de tor-dimension finie & I & 5.2\\
tor-indépendance & III & 1.5\\
trace & I & 8.1\\
de \(\mathcal C_0\)-type fini & I & 1.2\\
variété mixte de classe \(C^r\) & II App. II & 1.1\\
\bottomrule
\end{longtable}
\endgroup
\newpage
\section*{Index des notations}
\begingroup\small
\begin{longtable}{@{}>{\raggedright\arraybackslash}p{0.50\linewidth}>{\raggedright\arraybackslash}p{0.17\linewidth}>{\raggedright\arraybackslash}p{0.15\linewidth}@{}}
\toprule
 & Exposé & §\\
\midrule
\endhead
\(A_N\) & V & 5.5\\
\(\ch\) & V & 6.3\\
\(\Ch(A),\ \Ch_K(A)\) & V & 6.2\\
\((\Ch(A))_n\) & V & 6.5\\
\(\cl_{\mathcal C}\) & IV & 1.1\\
\(\Coh(S)\) & II & 2.2.1\\
\(\mathcal C_0\text{-}\parf\text{-}\amp(F)\) & I & 4.7\\
\(C^p_{\num}(X)\) & XIII & 5.2\\
\(C^r(V),\ C^\infty(V)\) & II App. II & 1.1\\
\(c(x),\ c^i(x),\ \widetilde c(x)\) & V & 6.7\\
\(\Dcat(A,S),\ \Dcat(A_S)\) & I & 2.15\\
\(\Dcat(A_{\coh}),\ \Dcat(A\text{-}n\text{-}\coh)\) & I & 2.15\\
\(\Dcat(A_S)_{\parf}\) & I & 4.9\\
\(\Dcat^{-}(A_S)_{\torf},\ \Dcat^{-}(A_X)_{\torf}\) & I & 5.3\\
\(\Dcat(\mathcal C)_{\mathcal C_0\text{-}\coh},\ \Dcat(\mathcal C)_{\mathcal C_0\text{-}n\text{-}\coh}\) & I & 2.6\\
\(\Dcat(\mathcal C)_{\mathcal C_0\text{-}\parf}\) & I & 4.9\\
\(\Dcat(f)_{\coh}\) & III & 1.2\\
\(\Dcat(f)_{n\text{-}\coh}\) & III & 1.2\\
\(\Dcat(f)_{\parf},\ \Dcat(f)_{\parf}\) & III & 4.2\\
\(\Dcat(f)_{\torf}\) & III & 3.2\\
\(\Diff^n(E,F)\) & II App. II & 1.2.2\\
\(d(L)\) & XIII & 2.2\\
\(d^n_{X/S}\) & II App. II & 1.2.1\\
\(d^n_{X/S}(E)\) & II App. II & 1.2.2\\
\(\Dcat^{-}(\Qcoh(S))_{\coh},\ \Dcat(\Qcoh(S))_{\parf}\) & II & 2.2.1\\
\(\Dcat(S),\ \Dcat(\underline S),\ \Dcat(S,A)\) & I & 2.15\\
\(\Dcat(X)_{f\text{-}\parf},\ \Dcat(X)_{\underline f\text{-}\parf}\) & III & 4.2\\
\(\Dcat(X)_{f\text{-}\torf}\) & III & 3.2\\
\(\Dcat(X)_{\qcoh}\) & III & 1.7\\
\(\Dcat(X)_{Y\text{-}\coh},\ \Dcat(X)_{Y\text{-}n\text{-}\coh}\) & III & 1.2\\
\(\Dcat(X)_{Y\text{-}\parf},\ \Dcat(X)_{Y\text{-}\parf}\) & III & 4.2\\
\(\Dcat(X)_{Y\text{-}\torf}\) & III & 3.2\\
\(\Dcat(Y)_{n\text{-}\coh}\) & III & 1.1\\
\(E^\vee\) & I & 7.3\\
\(*\) & V & 6.1\\
\(F|\) & X & 4.3.1\\
\(f^*\) (entre les \(K^\bullet\)) & IV & 2.7\\
\(f^*\) (entre les \(K_\bullet\)) & IV & 2.12\\
\(f_*\) (entre les \(K_\bullet\)) & IV & 2.11\\
\(f_*\) (entre les \(K^\bullet\)) & IV & 2.12\\
\(f_*\) (cas du fibré projectif) & VI & 5.2\\
\(\operatorname{gr}\) & VIII & 3.2\\
\(\Fil_n(X)\) & X & 1.1.1\\
\(F^{\prime >n},\ F'_{\leq n}\) & I & 2.12\\
\(\varphi^a\) & V & 2.9\\
\(\circ,\ \widehat G_K\) & V & 1.1\\
\(\circ,\ \widehat G_{\nabla}\) & V & 3.7\\
\(a\) & V & 1.3.3\\
\(g\) & V & 1.3.2\\
\(\operatorname{Gr}^\bullet(X),\ \operatorname{Gr}_\bullet(X)\) & X & 3.1\\
\(\operatorname{Gr}^\bullet_{\alg},\ \operatorname{Gr}^{\alg}_\bullet\) & X App. & 7.11\\
\(\operatorname{Gr}^\bullet_{\num},\ \operatorname{Gr}^{\num}_\bullet\) & X App. & 7.9.4\\
\(\gamma^n,\ \gamma_t\) & V & 3.2\\
\(\gamma^p(N,x)\) & V & 5.5\\
\(\widehat G^+,\ \widehat G^u\) & V & 2.3\\
\(\underline{\Hom}_{\Lambda}(E,F)\) & I & 4.5\\
\(k(A)\) & IV & 1.5\\
\(K^\bullet_{\alg},\ K^{\alg}_\bullet\) & X App. & 7.11\\
\(K^\bullet(A_X),\ K_\bullet(A_X)\) & IV & 2.2\\
\(K^\bullet(A_X)_{\naif},\ K_\bullet(A_X)_{\naif}\) & IV & 2.2\\
\(k(C)\) & I & 6.3\\
\(K^\bullet(f),\ K_\bullet(f)\) & IV & 3.3.1\\
\(K^\bullet_{\num},\ K^{\num}_\bullet\) & X App. & 7.9.4\\
\(K^\bullet(X),\ K_\bullet(X)\) & IV & 2.2\\
\(K^\bullet(X)_{\naif},\ K_\bullet(X)_{\naif}\) & IV & 2.2\\
\(K^\bullet(X/Y),\ K_\bullet(X/Y)\) & IV & 3.3.1\\
\(\Lib(U)\) & II & 2.0\\
\(\Loclib(U)\) & II & 2.0\\
\(L^{X/Y}\) & VIII & 2.1\\
\(\lambda^i\) & V & 2.1\\
\(\lambda_{-1}(N)\) & V & 5.1\\
\(\lambda_{-1}(N)\) & VII & 2.7\\
\(\lambda^p(N,x)\) & V & 5.3\\
\(\lambda_t\) & V & 2.1\\
\(\Mod(U)\) & II & 2.0\\
\(\parfamp_f(E),\ \parfamp_Y(E)\) & III & 4.1\\
\(\Parf(A_S),\ \Parf(\mathcal C,\mathcal C_0),\ \Parf(S)\) & I & 4.9\\
\(\Pic^N(X)\) & X & 4.5.2\\
\(\Pic^\circ(X)\) & X & 2.4.1\\
\(\Pic^\circ_{X/k}\) & X & 2.4\\
\(\Pic^\tau(X)\) & X & 2.4.1\\
\(\Pic(X)\) & X & 2.4\\
\(\underline{\Pic}_{X/k}\) & X & 2.4\\
\(P^n_{X/S}\) & II App. II & 1.2.1\\
\(P^n_{X/S}(E)\) & II App. II & 1.2.2\\
\(\Proj\) & VI & 2.3.2\\
\(\psi_k(x)\) & V & 7.1\\
\(\Qcoh(S)\) & II & 2.2.1\\
\(r^0_{\mathcal O_{X\times V}}\) & II App. II & 1.1\\
\(\rho_t\) & V & 3.4\\
\(\circ\) & V & 2.3\\
\(T_f\) & VIII & 2.5\\
\(\Todd\) & V & 1.16\\
\(\toramp(F)\) & I & 5.2\\
\(\toramp_f(E),\ \toramp_Y(E)\) & III & 3.1\\
\(\Tr_{E|U}\) & I & 8.1\\
\(T_{X/S}\) & II App. II & 1.2.1\\
\(\tau\) & I & 8.1\\
\(\nabla\) & V & 3.5\\
\(1+A[[t]]^+\) & V & 1.1\\
\(1+B^+\) & V & 1.11\\
\(X(n)\) & X & 1.1.2\\
\(x_N\) & V & 5.5\\
\(\langle x,y\rangle\) & X & 4.1\\
\(\mathbb Z(x)\) & X & 1.1.2\\
\bottomrule
\end{longtable}
\endgroup
\end{document}
